\section{作品概述}
\subsection{功能与特性}
系统完成$1280\times720$、30 fps相机采集、彩色重建、白平衡、自动曝光及七模式图像处理。主板将全帧图像缩放为RGB160输入，通过量化YOLO风格网络检测白细胞位置并判定五类，在视频中叠加目标框和英文类别名称。识别支路每两视频帧推理一次，中间帧保持完整结果。图像处理、检测分类、跨帧确认和标注均在FPGA完成，处理后的视频分别经网口送至PC与HDMI显示子板。

\subsection{应用领域}
作品主要面向血液细胞检测与形态学辅助分析，可用于血涂片图像的实时采集、观察、白细胞定位和分类。系统支持嗜碱性粒细胞、嗜酸性粒细胞、淋巴细胞、单核细胞和中性粒细胞五类识别，并将位置与类别直接标注在画面中，便于观察细胞形态、核对识别结果和保存典型图像。

针对血涂片观察中的照明与色彩差异，系统通过自动曝光和白平衡调节图像，并提供PC预览与HDMI显示两种观察方式，可作为血检辅助分析、血细胞形态学教学和相关检测算法验证的实验平台。

\subsection{主要技术特点}
系统在FPGA中构建流式ISP通路，完成RAW图像展开、Bayer彩色重建、白平衡、数字增益和多模式图像处理。自动白平衡根据校正前的色彩统计估计色偏，自动曝光根据校正后的亮度反馈调节曝光与增益，使图像调节与后续识别使用一致的处理链路。参数在帧边界统一生效，便于在连续视频中调整成像效果。

在模型部署方面，系统采用面向FPGA资源设计的轻量级检测分类网络，通过低位宽量化、流式首层计算和后续卷积运算单元复用，减少权重、特征存储及计算资源开销。模型直接使用ISP输出、未经风格化和文字叠加的图像，在板端完成白细胞定位、五类识别及结果确认；框与类别名称在视频发送前叠加，PC和HDMI端接收同一处理结果，无需PC参与实时推理。

\subsection{主要性能指标}
\begin{center}\small
\begin{tabularx}{\linewidth}{|L{2.5cm}|Y|}
\hline\tablehead 项目 & 指标\\\hline
图像输入 / 输出 & $1280\times720$，RAW10输入、RGB888输出，30 fps配置\\\hline
传输 / 显示 & 千兆UDP，1920包/帧；子板HDMI 720p60\\\hline
检测 / 分类 & 全帧$160\times160$ RGB输入；白细胞定位与五类识别\\\hline
推理 / 标注 & 15次/秒调度，非推理帧保持结果；目标框与英文类别名称\\\hline
模型量化 & 首尾层8位权重，中间层4位权重及激活，输入8位RGB\\\hline
整机资源 & Slice19016/21216，ERAM94/108，DSP35/40\\\hline
\end{tabularx}
\end{center}

\subsection{主要创新点}
\begin{enumerate}[label=（\arabic*）]
\item 国产FPGA内完成流式ISP与轻量级模型部署，实现白细胞检测分类和双端标注显示。
\item 抽帧推理与完整结果保持结合，按目标身份连续确认类别。
\item 统一帧身份贯穿自适应ISP、网络和显示，参数按帧提交。
\item 从数据清理、整数对拍到物理实现，以哈希关联模型与位流。
\end{enumerate}

\subsection{设计流程}
先建立采集、ISP、网络和显示链路，再接入全帧量化检测分类。WBCAtt+单一来源实验因实拍空场景误检撤回，随后使用WBCAtt+与BCCD两来源训练、量化和导出，完成整数对拍、720p调度检查、回归及物理实现后下载主板SRAM。当前工程统一收于\path{camera processing unit/}，主板、子板、训练和文档分目录管理；目录整理只调整工程引用路径，实时算法沿用最终候选实现。用户已反馈实拍改善并自行完成Flash烧录，独立定量相机验收与冷启动核验仍待完成。
